%=======================================================================
%  ENGINEERING ARENA 2026 -- FULL COMPETITION SOLUTION SET
%  Team: flop_flip
%  Single-file LaTeX deliverable. Compile with pdfLaTeX on Overleaf.
%=======================================================================
\documentclass[11pt]{article}

\usepackage[a4paper,margin=2.2cm]{geometry}
\usepackage{amsmath,amssymb}
\usepackage{xcolor}
\usepackage{graphicx}
\usepackage{enumitem}
\usepackage[most]{tcolorbox}
\usepackage{booktabs}
\usepackage{longtable}
\usepackage{array}
\usepackage[hidelinks]{hyperref}

\setlength{\parindent}{0pt}
\setlength{\parskip}{5pt}

% ---- Answer box: the exact input/action for the interactive window ----
\newtcolorbox{answer}{
  colback=yellow!12, colframe=orange!75!black,
  boxrule=1pt, arc=2pt, left=8pt, right=8pt, top=5pt, bottom=5pt,
  fonttitle=\bfseries, coltitle=white, colbacktitle=orange!75!black,
  title={INPUT TO ENTER / ACTION IN INTERACTIVE WINDOW}, breakable}

% ---- Info box: governing relations ----
\newtcolorbox{relbox}{
  colback=blue!5, colframe=blue!50!black,
  boxrule=0.6pt, arc=2pt, left=8pt, right=8pt, top=4pt, bottom=4pt,
  fonttitle=\bfseries, title={Governing relations}, breakable}

\newcommand{\task}[2]{\section{Task #1 --- #2}}

\title{\textbf{Engineering Arena 2026}\\[4pt]
\large Complete Solution Set --- Helping Stemo Rescue Athena}
\author{Team: \texttt{flop\_flip}}
\date{\today}

\begin{document}
\maketitle

\begin{abstract}
\noindent This document contains the worked solutions, governing equations and
calculations for all 18 interactive tasks of Engineering Arena 2026. For each
task the engineering reasoning is shown, followed by an \textbf{orange box}
giving the exact value to type or the action to perform in the interactive
web page. Tasks whose on-screen data is hidden (e.g.\ an offline scanner) are
solved symbolically, with the procedure to apply once the values are revealed.
\end{abstract}

\tableofcontents
\newpage

%=======================================================================
\task{1}{Column Stability --- Euler Buckling of a Braced Column}
%=======================================================================
\textbf{Problem.} A pinned escape-hatch column of length $L=4.0$\,m and
flexural rigidity $EI=80$\,kN$\cdot$m$^2$ carries an axial load
$P=56.0$\,kN. A movable brace at height $a$ splits the column into two
segments. The critical load is governed by the \emph{longer} unsupported
segment. We must position the brace so the column does not buckle.

\begin{relbox}
\[
P_{cr}=\frac{\pi^2 EI}{L_e^{2}},\qquad
L_e=\max(a,\;L-a),\qquad
\text{Win: } P_{cr}\ge P_{applied}.
\]
\end{relbox}

\textbf{Solution.} Rearranging the buckling condition for the maximum
admissible effective length:
\[
P_{cr}\ge P \;\Longrightarrow\;
L_e\le \pi\sqrt{\frac{EI}{P}}
      =\pi\sqrt{\frac{80}{56}}=3.755\ \text{m}.
\]
Since $L_e=\max(a,L-a)$, the brace must satisfy
$a\le 3.755$ and $L-a\le 3.755$, i.e. $0.245\le a\le 3.755$\,m.
The \emph{strongest} configuration places the brace at mid-height,
$a=L/2=2.0$\,m, which gives $L_e=2.0$\,m and
\[
P_{cr}=\frac{\pi^2(80)}{2.0^{2}}=197.4\ \text{kN},
\qquad
\text{SF}=\frac{197.4}{56}=3.53 .
\]
This is the ``brace at midpoint $\rightarrow P_{cr}\times 4$'' key insight:
halving $L_e$ quadruples $P_{cr}$.

\begin{answer}
Drag the brace to mid-height \textbf{$a = 2.00$\,m} (midpoint).\\
Result: $L_e = 2.00$\,m, $P_{cr}=197.4$\,kN $\ge 56$\,kN. Press \textbf{LOCK BRACE}.\\
(Any $a$ in $0.245$--$3.755$\,m satisfies the win condition; $a=2.0$\,m maximises $P_{cr}$.)
\end{answer}

%=======================================================================
\task{2}{Solar Sail MPPT --- Maximum Power Point Tracking}
%=======================================================================
\textbf{Problem.} A research-satellite solar panel operates on a non-linear
$P$--$V$ curve. The Maximum Power Point Tracker must sit at the voltage that
maximises output power. Sail $\alpha$: $I_{sc}=8$\,A, $V_{oc}=20$\,V,
irradiance $1000$\,W/m$^2$.

\begin{relbox}
\[
P(V)=V\cdot I(V),\qquad
\text{MPP: } \left.\frac{dP}{dV}\right|_{V=V_{mp}}=0 .
\]
$I_{sc}\propto$ irradiance; $V_{oc}$ varies weakly (logarithmically) with
irradiance and temperature.
\end{relbox}

\textbf{Solution.} Using a single-diode model $I(V)=I_{sc}\!\left[1-e^{(V-V_{oc})/V_t}\right]$,
the condition $dP/dV=0$ gives
\[
1=e^{(V_{mp}-V_{oc})/V_t}\Bigl(1+\tfrac{V_{mp}}{V_t}\Bigr).
\]
For sail $\alpha$ this is satisfied at $V_{mp}\approx16.26$\,V, where
$I\approx7.33$\,A and
\[
P_{max}=V_{mp}\,I(V_{mp})\approx16.26\times7.33\approx119.3\ \text{W}.
\]
This matches the engineering rule of thumb $V_{mp}\approx0.81\,V_{oc}$.
For the remaining sails the peak is found the same way; $I_{sc}$ scales with
irradiance ($\beta$: $420$\,W/m$^2\!\Rightarrow I_{sc}\approx3.36$\,A;
$\gamma$: $1400$\,W/m$^2\!\Rightarrow I_{sc}\approx11.2$\,A), so each curve's
peak shifts but stays near $V\approx0.8\,V_{oc}$.

\begin{answer}
Drag the tracker dot to the \textbf{peak of the gold $P$--$V$ curve}:\\
Sail $\alpha$: \textbf{$V_{mp}\approx16.26$\,V}, $P_{max}\approx119.3$\,W (100\,\% of $P_{max}$). Press \textbf{LOCK MPP \& HARVEST}.\\
For each subsequent sail, drag to that curve's peak ($V_{mp}\approx0.8\,V_{oc}$).
\end{answer}

%=======================================================================
\task{3}{Mohr's Circle Audit --- Principal Stresses}
%=======================================================================
\textbf{Problem.} A stress element with components $(\sigma_x,\sigma_y,\tau_{xy})$
must be analysed by constructing Mohr's circle to expose the principal
stresses $\sigma_1,\sigma_2$. \emph{(In the screenshot the structural scanner
is offline; the procedure below applies once \textbf{SCAN} reveals the element
values.)}

\begin{relbox}
\[
C=\frac{\sigma_x+\sigma_y}{2},\quad
R=\sqrt{\left(\frac{\sigma_x-\sigma_y}{2}\right)^{2}+\tau_{xy}^{2}},
\]
\[
\sigma_1=C+R,\qquad \sigma_2=C-R,\qquad \tau_{max}=R .
\]
\end{relbox}

\textbf{Solution (worked template).} After pressing SCAN, read
$\sigma_x,\sigma_y,\tau_{xy}$ from the top-right stress element. For example,
with the representative values $\sigma_x=120$\,MPa, $\sigma_y=40$\,MPa,
$\tau_{xy}=30$\,MPa:
\[
C=\frac{120+40}{2}=80\ \text{MPa},\qquad
R=\sqrt{40^{2}+30^{2}}=50\ \text{MPa},
\]
\[
\sigma_1=80+50=130\ \text{MPa},\qquad
\sigma_2=80-50=30\ \text{MPa},\qquad \tau_{max}=50\ \text{MPa}.
\]
The circle is centred on the $\sigma$-axis at $(C,0)$ with radius $R$; it must
pass through both diamond points $A(\sigma_x,+\tau_{xy})$ and
$B(\sigma_y,-\tau_{xy})$.

\begin{answer}
1. Drag the resize handle until the orange circle radius equals
$R=\sqrt{\left(\frac{\sigma_x-\sigma_y}{2}\right)^2+\tau_{xy}^2}$
(circle passes through both blue/purple diamonds).\\
2. Click the \textbf{$\sigma_1$ node} (right intercept, $\sigma_1=C+R$) and the
\textbf{$\sigma_2$ node} (left intercept, $\sigma_2=C-R$).\\
3. Press \textbf{COMPLETE CHECKLIST}.\\
\textit{Template example} ($\sigma_x{=}120,\sigma_y{=}40,\tau_{xy}{=}30$):
$R=50$, $\sigma_1=130$\,MPa, $\sigma_2=30$\,MPa.
\end{answer}

%=======================================================================
\task{4}{Radar Convolution --- Cross-Correlation Peak}
%=======================================================================
\textbf{Problem.} A radar echo $x(t)$ contains three blips; only one is the
real return (the others are decoy pods). Sliding a search pulse $h(t)$ over the
echo, the real target produces the \emph{global maximum} of the
cross-correlation.

\begin{relbox}
\[
R(\tau)=\sum_t h(t)\,x(t+\tau),\qquad
\tau^{*}=\operatorname*{arg\,max}_{\tau} R(\tau).
\]
$h(t)=$ search pulse, $x(t)=$ radar echo.
\end{relbox}

\textbf{Solution.} This is a matched-filter detection problem. Cross-correlation
measures similarity between the pulse and the echo at every lag $\tau$. The
\emph{real} wizard return has the same shape as the search pulse, so the
correlation reaches its highest peak there; decoys, being only partially
similar, yield smaller secondary peaks. The optimal lag $\tau^{*}$ is the
abscissa of the tallest peak of $R(\tau)$ --- on the cross-correlation panel
this is the dominant central peak, clearly higher than the two flanking
(decoy) peaks.

\begin{answer}
Drag the amber search pulse so the cyan correlation dot climbs to the
\textbf{tallest (global-maximum) peak} of $R(\tau)$ --- the dominant central
lobe, not the lower side peaks. Stop within $\pm10$ samples of that peak
(dot turns green) and press \textbf{LOCK TARGET}.
\end{answer}

%=======================================================================
\task{5}{Arcane Capacitor --- Displacement-Current B-Field}
%=======================================================================
\textbf{Problem.} A time-varying $E$-field in a capacitor gap produces a
displacement current and hence a magnetic field. Tune the AC frequency so the
rim field $B_{rim}$ matches the target.\\
Live data: $f=69.0$\,MHz, $B_{rim}=127{,}952.0$\,nT,
$B_{target}=195{,}896.8$\,nT, error $=34.7\,\%$.

\begin{relbox}
\[
I_d=\varepsilon_0 A\,(2\pi f)\,E_0,\qquad
B=\frac{\mu_0 I_d}{2\pi r}\;\;\Longrightarrow\;\; B\propto f .
\]
\end{relbox}

\textbf{Solution.} Because every other quantity is fixed, the rim field is
\emph{linearly proportional to frequency}, $B=k f$ with
$k=B_{rim}/f=127{,}952.0/69.0=1854.4$\,nT/MHz. The frequency that hits the
target is
\[
f_{target}=f\cdot\frac{B_{target}}{B_{rim}}
=69.0\times\frac{195{,}896.8}{127{,}952.0}
=69.0\times1.5310=105.6\ \text{MHz}.
\]
This lies well inside the $\pm5\,\%$ window
($B$ between $186.1$ and $205.7\,\mu$T $\Rightarrow f\in[100.4,110.9]$\,MHz).

\begin{answer}
Drag the frequency slider to \textbf{$f \approx 105.6$\,MHz}\\
(acceptable range $100.4$--$110.9$\,MHz for $\pm5\,\%$). Press \textbf{LOCK ON}.
\end{answer}

%=======================================================================
\task{6}{Counter-Flow Heat Exchanger --- NTU-$\varepsilon$ Method}
%=======================================================================
\textbf{Problem.} Cool the reactor coolant outlet below $75\,^{\circ}$C.
Inlets: $T_{h,in}=200\,^{\circ}$C, $T_{c,in}=20\,^{\circ}$C;
$UA=5$\,kW/K, $c_p=4.18$\,kJ/(kg$\cdot$K). Adjust the hot- and cold-side
mass flow rates $\dot m_h,\dot m_c$.

\begin{relbox}
\[
C=\dot m c_p,\qquad \text{NTU}=\frac{UA}{C_{min}},\qquad C_r=\frac{C_{min}}{C_{max}},
\]
\[
\varepsilon=\frac{1-e^{-\text{NTU}(1-C_r)}}{1-C_r\,e^{-\text{NTU}(1-C_r)}}
\quad(\text{counter-flow}),
\]
\[
Q=\varepsilon\,C_{min}\,(T_{h,in}-T_{c,in}),\qquad
T_{h,out}=T_{h,in}-\frac{Q}{C_h}.
\]
\end{relbox}

\textbf{Solution.} Lowering $\dot m_h$ reduces $C_h$: this both raises NTU
($=UA/C_{min}$) and lowers $C_r$, and both effects increase $\varepsilon$ and
the hot-side temperature drop. Raising $\dot m_c$ further lowers $C_r$.
With the hot side as $C_{min}$, $T_{h,out}=T_{h,in}-\varepsilon\,\Delta T_{max}$,
so the win needs $\varepsilon\ge125/180=0.694$.

Taking $\dot m_h=0.80$\,kg/s and $\dot m_c=2.50$\,kg/s:
\[
C_h=3.34,\;\; C_c=10.45\ \text{kW/K},\quad C_r=0.320,\quad
\text{NTU}=\frac{5}{3.34}=1.50,
\]
\[
\varepsilon=\frac{1-e^{-1.50(0.68)}}{1-0.320\,e^{-1.50(0.68)}}=0.722,\qquad
Q=0.722\times3.34\times180=434\ \text{kW},
\]
\[
T_{h,out}=200-\frac{434}{3.34}=70.1\,^{\circ}\text{C}\;<\;75\,^{\circ}\text{C}.\;\checkmark
\]
(The exact boundary $T_{h,out}=75\,^{\circ}$C occurs at
$\dot m_h\approx0.86$\,kg/s; $0.80$\,kg/s gives a safe margin.)

\begin{answer}
Set the hot-flow slider to \textbf{$\dot m_h \approx 0.80$\,kg/s} (toward minimum)
and the cold-flow slider to \textbf{$\dot m_c \approx 2.50$\,kg/s} (toward maximum).\\
Result: $\varepsilon\approx72\,\%$, $T_{h,out}\approx70\,^{\circ}$C. Press \textbf{TRIGGER SHUTDOWN}.
\end{answer}

%=======================================================================
\task{7}{Sentry-Bot MK-VII --- 3-Joint Forward Kinematics}
%=======================================================================
\textbf{Problem.} Place the end-effector of a 3-joint arm
($\theta_1$ yaw, $\theta_2$ shoulder, $\theta_3$ elbow) within 12 units of the
Obsidian Lock at $(x,y,z)=(167.01,\,50.41,\,-16.23)$.

\begin{relbox}
\[
x=\cos\theta_1\!\cdot\! r,\quad y=\sin\theta_1\!\cdot\! r,\quad
r=L_1\cos\theta_2+L_2\cos(\theta_2+\theta_3),
\]
\[
z=L_1\sin\theta_2+L_2\sin(\theta_2+\theta_3).
\]
\end{relbox}

\textbf{Solution.} \emph{Calibration.} The displayed pose
$\theta_1{=}16.6^{\circ},\theta_2{=}-3.1^{\circ},\theta_3{=}-4.5^{\circ}$ yields
$EE=(186.0,55.5,-17.3)$. From the planar pair of equations this fixes the link
lengths $L_1\approx108.6$, $L_2\approx86.4$ units.

\emph{Inverse kinematics for the target.}
\[
\theta_1=\operatorname{atan2}(50.41,167.01)=16.8^{\circ},
\]
\[
r_t=\sqrt{167.01^2+50.41^2}=174.5,\quad
R_t=\sqrt{r_t^2+z_t^2}=175.2,
\]
\[
\cos\theta_3=\frac{R_t^2-L_1^2-L_2^2}{2L_1L_2}=0.609
\;\Rightarrow\;\theta_3=-52.5^{\circ},
\]
\[
\theta_2=\operatorname{atan2}(z_t,r_t)
+\arccos\!\frac{L_1^2+R_t^2-L_2^2}{2L_1R_t}=17.7^{\circ}.
\]
Forward-checking these angles returns $EE\approx(167.0,50.4,-16.2)$ ---
essentially on target (dist $\approx0.1\ll12$).

\begin{answer}
Set the canvas sliders to:\\
\textbf{$\theta_1\approx16.8^{\circ}$ (yaw)},\quad
\textbf{$\theta_2\approx17.7^{\circ}$ (shoulder)},\quad
\textbf{$\theta_3\approx-52.5^{\circ}$ (elbow)}.\\
The elbow must fold in sharply (reach drops $195\to175$). Press \textbf{UNLOCK DOOR}.
\end{answer}

%=======================================================================
\task{8}{Rankine Cycle --- Thermal Efficiency}
%=======================================================================
\textbf{Problem.} Engineer a Rankine cycle exceeding $32\,\%$ thermal
efficiency by adjusting boiler temperature $T_H$ and condenser temperature
$T_L$. Start: $T_H=180\,^{\circ}$C, $T_L=90\,^{\circ}$C, $\eta=18\,\%$.

\begin{relbox}
\[
\eta=\frac{W_{net}}{Q_{in}},\qquad
\eta_{Carnot}=1-\frac{T_L}{T_H}\quad(\text{absolute temperatures}).
\]
\end{relbox}

\textbf{Solution.} Raising $T_H$ and lowering $T_L$ enlarges the enclosed area
of the $T$--$s$ cycle, which is the net work $W_{net}$, while the Carnot
ceiling rises. The current cycle ($\eta=18\,\%$) sits close to its Carnot bound
$1-363/453=19.9\,\%$. Driving the bars to their extremes
$T_H=350\,^{\circ}$C ($623$\,K) and $T_L=30\,^{\circ}$C ($303$\,K):
\[
\eta_{Carnot}=1-\frac{303}{623}=51.4\,\%.
\]
Maintaining the same fraction of the Carnot limit
($\eta/\eta_{Carnot}\approx0.9$) gives an actual cycle efficiency
$\eta\approx0.9\times51.4\approx46\,\%$ --- comfortably above the $32\,\%$
target.

\begin{answer}
Drag the \textbf{red $T_H$ bar up to $\approx350\,^{\circ}$C} and the
\textbf{blue $T_L$ bar down to $\approx30\,^{\circ}$C}.\\
Result: $\eta_{Carnot}=51.4\,\%$, cycle $\eta\approx45$--$46\,\% > 32\,\%$.
Press \textbf{ACTIVATE OVERDRIVE}.
\end{answer}

%=======================================================================
\task{9}{Nyquist Hack --- Encirclement Count}
%=======================================================================
\textbf{Problem.} For Shield Alpha with open-loop transfer function
$G(s)=\dfrac{1}{(s+1)^2}$, count the clockwise encirclements $N$ of the
critical point $(-1,0)$ and determine stability.

\begin{relbox}
\[
Z=N+P,\qquad
\begin{cases}
N=\#\text{ clockwise encirclements of }(-1,0)\\
P=\#\text{ open-loop unstable poles}\\
Z>0 \Rightarrow \text{UNSTABLE}
\end{cases}
\]
\end{relbox}

\textbf{Solution.} $G(s)=1/(s+1)^2$ has a double pole at $s=-1$, in the
left-half plane, so $P=0$. The Nyquist locus starts at
$G(j0)=1\,(=1\angle0^{\circ})$ and spirals into the origin as
$\omega\to\infty$, with phase sweeping only from $0^{\circ}$ to $-180^{\circ}$.
The plotted curve is a small loop confined near the origin and the point
$(+1,0)$; it never wraps around the critical point $(-1,0)$. Hence
\[
N=0,\qquad Z=N+P=0\quad\Rightarrow\quad\text{closed-loop STABLE}.
\]

\begin{answer}
Select \textbf{$N = 0$}.\\
(The locus does not encircle $(-1,0)$; with $P=0$, $Z=N+P=0$ --- Shield Alpha is stable.)
\end{answer}

%=======================================================================
\task{10}{Torsional Stress --- Yield Torque of a Shaft}
%=======================================================================
\textbf{Problem.} Find the exact torque $T_{max}$ that brings a solid vault
shaft to its yield point at the outer fibre.\\
Data: $R=53$\,mm, $J=12.394\times10^{-6}$\,m$^4$,
$\tau_{yield}=130$\,MPa.

\begin{relbox}
\[
\tau(r)=\frac{T\,r}{J},\qquad J=\frac{\pi R^4}{2},\qquad
\tau_{max}=\frac{T R}{J}\;\Rightarrow\;
T_{max}=\frac{\tau_{yield}\,J}{R}.
\]
\end{relbox}

\textbf{Solution.} Shear stress varies linearly with radius (zero at the
centre, maximum at the surface); only the outer fibre reaches yield first.
Setting $\tau_{max}=\tau_{yield}$:
\[
T_{max}=\frac{\tau_{yield}\,J}{R}
=\frac{(130\times10^{6})(12.394\times10^{-6})}{0.053}
=30{,}400\ \text{N}\cdot\text{m}\approx30.40\ \text{kN}\cdot\text{m}.
\]
\emph{Check:} at $T=29.992$\,kN$\cdot$m the model reads
$\tau_{max}=128$\,MPa $=99\,\%$ of $\tau_{yield}$ --- consistent.

\begin{answer}
Drag the torque slider to \textbf{$T_{max} = 30.40$\,kN$\cdot$m}
($\approx30{,}400$\,N$\cdot$m).\\
The green dot meets the red yield line at $r=R$ ($\tau_{max}=130$\,MPa).
Press \textbf{APPLY T\_MAX}.
\end{answer}

%=======================================================================
\task{11}{Void Frequency Tuner --- RLC Resonance}
%=======================================================================
\textbf{Problem.} An RLC filter must be tuned by sliding the inductor knob $L$
so the receiver's resonant frequency $f_0$ aligns with the wizard's spike
$f_T$. Data: $C=1.0\,\mu$F, $Q=12$, current $L=3.57$\,mH giving
$f_0=2663$\,Hz at $100\,\%$ signal.

\begin{relbox}
\[
f_0=\frac{1}{2\pi\sqrt{LC}}\qquad\Longrightarrow\qquad
L=\frac{1}{(2\pi f_T)^{2}\,C}.
\]
\end{relbox}

\textbf{Solution.} Resonance occurs where inductive and capacitive reactances
cancel. To shift the green resonant peak onto the red target spike, set $L$
from the inverted resonance formula. At the displayed setting
\[
f_0=\frac{1}{2\pi\sqrt{(3.57\times10^{-3})(1.0\times10^{-6})}}=2664\ \text{Hz},
\]
which already coincides with $f_T$ --- the receiver reads SIGNAL $100\,\%$,
so $f_T\approx2663$\,Hz and $L=3.57$\,mH is the correct inductance.

\begin{answer}
Set the inductor knob to \textbf{$L = 3.57$\,mH} $\Rightarrow f_0=2663$\,Hz $=f_T$
(Signal $100\,\%$). Press \textbf{LOCK SIGNAL}.\\
General retune: $L=\dfrac{1}{(2\pi f_T)^2 C}$.
\end{answer}

%=======================================================================
\task{12}{Beam Deflection --- Location of Maximum Deflection}
%=======================================================================
\textbf{Problem.} A simply supported bridge beam ($L=6.00$\,m,
$R_a=29.5$\,kN, $R_b=42.1$\,kN) carries three point loads. Predict the
position of maximum deflection (tolerance $\pm0.55$\,m).

\begin{relbox}
\[
EI\,y''=M(x),\qquad
\delta_{max}\ \text{where}\ y'(x)=0,
\]
$V(x)=0\Rightarrow dM/dx=0\Rightarrow M(x)$ is maximum.
\end{relbox}

\textbf{Solution.} For a simply supported beam the elastic curve is flattest
(zero slope, maximum deflection) very close to the section of maximum bending
moment. The bending-moment diagram peaks at $x\approx3.06$\,m
(where the shear $V(x)$ changes sign under the dominant $31$\,kN load).
Hence the maximum deflection occurs at
\[
x_{\delta_{max}}\approx3.06\ \text{m},
\]
essentially mid-span ($L/2=3.00$\,m), comfortably within the $\pm0.55$\,m
tolerance band.

\begin{answer}
Drag the amber marker to \textbf{$x \approx 3.06$\,m}
(the BMD peak; $3.00$\,m mid-span is also within tolerance).
Press \textbf{REVEAL DEFLECTION}.
\end{answer}

%=======================================================================
\task{13}{Singularity Trap --- Jacobian Determinant}
%=======================================================================
\textbf{Problem.} A 2-link arm must reach the gold target while keeping the
normalised Jacobian determinant above threshold to avoid a kinematic
singularity.

\begin{relbox}
\[
\det(\mathbf J)=L_1 L_2\sin\theta_2,\qquad
\frac{|\det(\mathbf J)|}{L_1 L_2}=|\sin\theta_2|>0.25 .
\]
$\theta_2=0^{\circ}$ (fully extended) or $\pm180^{\circ}$ (fully folded) $\Rightarrow$ singularity.
\end{relbox}

\textbf{Solution.} The Jacobian determinant collapses to zero when the arm is
straight or fully folded, where the manipulator loses a degree of freedom.
The safety condition $|\sin\theta_2|>0.25$ requires
\[
14.5^{\circ}<\theta_2<165.5^{\circ}.
\]
Manipulability is maximised at $\theta_2=90^{\circ}$ ($|\det|/L_1L_2=1$).
Reaching the target with the elbow bent near $\theta_2\approx87^{\circ}$
(as in the displayed pose, $\theta_1=8.5^{\circ}$) keeps
$|\det|/L_1L_2\approx0.998$ --- far from any singularity zone.

\begin{answer}
Drag the Elbow/EE handles so the arm reaches the gold target with the elbow
well-bent: \textbf{$\theta_2\approx87^{\circ}$} (keep $14.5^{\circ}<\theta_2<165.5^{\circ}$;
$90^{\circ}$ is ideal), $\theta_1\approx8.5^{\circ}$.\\
Result: dist $<28$\,px and $|\det \mathbf J|/L_1L_2\approx1$. Press \textbf{ESCAPE!}.
\end{answer}

%=======================================================================
\task{14}{PID Stabilizer --- Tuning the Wind-Void Crossing}
%=======================================================================
\textbf{Problem.} Tune the PID gains $K_p,K_i,K_d$ to reduce the oscillation
energy below $1250$. Start: $K_p=0.50$, $K_i=0.10$, $K_d=0.30$, $E=1274.1$.

\begin{relbox}
\[
u(t)=K_p\,e(t)+K_i\!\int e(t)\,dt+K_d\,\frac{de}{dt}.
\]
$K_p$: reacts to current error; $K_i$: removes accumulated drift;
$K_d$: damps oscillation by anticipating trends.
\end{relbox}

\textbf{Solution.} The flight path shows a lightly damped oscillation, and the
energy ($E=1274$) is only marginally above the $1250$ target. The derivative
term is the one that suppresses oscillation, so it should be increased; the
proportional term is slightly too aggressive (it feeds the overshoot) and is
trimmed; the integral term is kept small to avoid integral wind-up that would
re-introduce slow oscillation:
\[
K_d\uparrow,\qquad K_p\downarrow\ \text{slightly},\qquad K_i\downarrow.
\]
A balanced set near $K_p\approx0.40$, $K_i\approx0.05$, $K_d\approx0.55$
adds damping without making the craft sluggish, pulling $E$ below $1250$.

\begin{answer}
Set the dials to approximately \textbf{$K_p\approx0.40$,\quad
$K_i\approx0.05$,\quad $K_d\approx0.55$}\\
(increase $K_d$ to damp oscillation, ease $K_p$ down, keep $K_i$ small).
Fine-tune until OSCILLATION ENERGY $<1250$, then press \textbf{ATTEMPT CROSSING}.
\end{answer}

%=======================================================================
\task{15}{Mach Shock --- Normal Shock in a C-D Nozzle}
%=======================================================================
\textbf{Problem.} Position a normal shock in the diverging section of a
converging--diverging nozzle so the exit pressure ratio hits
$P_{exit}/P_0=0.738$ ($\pm4\,\%$).
Data: $\gamma=1.4$, exit area ratio $A_e/A^{*}=2.125$, shock at $x>0.5$.

\begin{relbox}
\[
\frac{A}{A^{*}}=\frac{1}{M}\!\left[\frac{2}{\gamma+1}\!\left(1+\tfrac{\gamma-1}{2}M^2\right)\right]^{\frac{\gamma+1}{2(\gamma-1)}},
\]
\[
M_2^2=\frac{M_1^2+5}{7M_1^2-1},\qquad
\frac{P_{0,2}}{P_{0,1}}\ \text{(normal-shock stagnation loss)}.
\]
\end{relbox}

\textbf{Solution.} Moving the shock downstream raises the pre-shock Mach number
$M_1$, strengthening the shock and increasing the stagnation-pressure loss; the
subsonic flow after the shock then expands to a different exit pressure.
Sweeping the shock area ratio $A_s/A^{*}$:
\[
\frac{A_s}{A^{*}}=1.40\Rightarrow\frac{P_{exit}}{P_0}=0.760,\qquad
\frac{A_s}{A^{*}}=1.50\Rightarrow0.715,
\]
so the target $0.738$ is reached at $A_s/A^{*}\approx1.445$
($M_1\approx1.81$). With the fitted nozzle geometry
$A/A^{*}=1+4.5\,(x-0.5)^2$, this corresponds to
\[
x=0.5+\sqrt{\frac{1.445-1}{4.5}}\approx0.81 .
\]
The full normal-shock calculation at this station returns
$P_{exit}/P_0\approx0.740$ --- inside the $\pm4\,\%$ window.

\begin{answer}
Drag the shock marker into the diverging section to \textbf{$x\approx0.81$}
(area ratio $A_s/A^{*}\approx1.45$).\\
Result: $P_{exit}/P_0\approx0.738$. Press \textbf{FIRE HACK}.
\end{answer}

%=======================================================================
\task{16}{Kalman Teleportation Tracker --- Optimal Estimate}
%=======================================================================
\textbf{Problem.} Each round, noisy yellow sensor blips appear; click where the
teleporting wizard truly is. The Kalman filter then fuses the estimate.
Parameters: $Q=8$ (process noise), $R=15$ (sensor noise), initial $P=50$.

\begin{relbox}
Predict: $\hat x_k^-=\hat x_{k-1}$, $P_k^-=P_{k-1}+Q$.\\
Update: $K_k=\dfrac{P_k^-}{P_k^-+R}$,\;\;
$\hat x_k=\hat x_k^-+K_k\,(z_k-\hat x_k^-)$,\;\;
$P_k=(1-K_k)P_k^-$.
\end{relbox}

\textbf{Solution.} On the first round the prior is highly uncertain
($P=50\gg R$), so the Kalman gain is large ($K\approx50/65\approx0.77$) and the
estimate is dominated by the measurements. With several independent blips the
optimal estimate is their \emph{average} (the maximum-likelihood location of
the source). Therefore click at the \textbf{centroid of the yellow blip
cluster}. As rounds progress, $P$ shrinks, $K$ falls, and the estimate
increasingly trusts the converged prior --- but the click rule (centre of the
blip pattern, nudged toward the running estimate) stays the same.

\begin{answer}
Click on the grid at the \textbf{centroid (average position) of the visible
yellow blips} --- for Round 1 the blips cluster near $x\approx33$--$38$, so
click at \textbf{$x\approx35$}.\\
Each round: click the blip-cluster centre; the filter converges ($P\downarrow$)
over the 7 rounds.
\end{answer}

%=======================================================================
\task{17}{Spectral De-Scrambler --- FFT Carrier Frequencies}
%=======================================================================
\textbf{Problem.} The raw audio is white noise in the time domain; switching to
the frequency domain (FFT) reveals the wizard's jamming spikes. Click the three
sharp resonant peaks to lock the carrier frequencies.

\begin{relbox}
\[
X(k)=\sum_{n} x(n)\,e^{-i\,2\pi kn/N}.
\]
$|X(k)|$ isolates the amplitude of periodic sine components buried in random
noise $x(n)$.
\end{relbox}

\textbf{Solution.} A Fourier transform decomposes the time signal into its
constituent frequencies. Random noise spreads energy thinly across all bins,
while a periodic carrier concentrates energy into a single sharp spike. The FFT
analyzer shows three such resonant peaks; reading their abscissae gives the
three hidden carrier frequencies:
\[
f_1\approx45\ \text{Hz},\qquad f_2\approx212\ \text{Hz},\qquad
f_3\approx395\ \text{Hz}.
\]

\begin{answer}
Click the three sharp peaks in the FFT panel at
\textbf{$f_1\approx45$\,Hz}, \textbf{$f_2\approx212$\,Hz},
\textbf{$f_3\approx395$\,Hz} (CH\_1, CH\_2, CH\_3). Then press \textbf{DECODE}.
\end{answer}

%=======================================================================
\task{18}{The Prism Shield --- LQR Controller Q/R Balance}
%=======================================================================
\textbf{Problem.} Select the LQR weighting on the state-space plot to balance
fast stabilisation against efficient energy use.

\begin{relbox}
\[
J=\int_0^{\infty}\!\bigl(\mathbf x^{\!\top}\mathbf Q\,\mathbf x
+\mathbf u^{\!\top}\mathbf R\,\mathbf u\bigr)\,dt,
\qquad \dot{\mathbf x}=A\mathbf x+B\mathbf u .
\]
$\mathbf Q$ penalises state error (speed); $\mathbf R$ penalises control effort (energy).
\end{relbox}

\textbf{Solution.} The LQR cost trades two competing objectives. A large
$\mathbf Q$ (relative to $\mathbf R$) drives the state to zero quickly but
demands aggressive, energy-hungry control; a large $\mathbf R$ saves energy but
stabilises slowly. The mission requires the platform to settle \emph{before the
shield rotates} yet still \emph{retain energy for the final blast} --- so
neither extreme is acceptable. The sweet spot is a roughly balanced ratio,
$\mathbf Q/\mathbf R\approx1$ (with a mild bias toward $\mathbf Q$ for prompt
settling). On the Position--Velocity selector this is a point near the centre,
slightly into the high-$\mathbf Q$ region, giving a well-damped trajectory that
spirals smoothly into the origin without overshoot.

\begin{answer}
Click the state-space plot at a \textbf{balanced point} --- moderate
$\mathbf Q$ and moderate $\mathbf R$ ($\mathbf Q/\mathbf R\approx1$,
slightly biased toward higher $\mathbf Q$): near the plot centre, just into the
high-$\mathbf Q$ side. Press \textbf{TEST CONFIG}.\\
(Avoid the extremes: pure high-$\mathbf Q$ = overshoot/energy waste;
pure high-$\mathbf R$ = too slow.)
\end{answer}

%=======================================================================
\section{Summary of All Inputs}
%=======================================================================
\renewcommand{\arraystretch}{1.35}
\begin{longtable}{@{}p{0.6cm}p{4.4cm}p{10.5cm}@{}}
\toprule
\textbf{\#} & \textbf{Task} & \textbf{Value to enter / action}\\
\midrule
\endhead
1  & Column Stability      & Brace at $a=2.00$\,m; $P_{cr}=197.4$\,kN.\\
2  & Solar Sail MPPT       & Tracker at $V_{mp}\approx16.26$\,V, $P\approx119.3$\,W (sail $\alpha$).\\
3  & Mohr's Circle         & $R=\sqrt{((\sigma_x-\sigma_y)/2)^2+\tau_{xy}^2}$; click $\sigma_1=C+R$, $\sigma_2=C-R$.\\
4  & Radar Convolution     & Drag pulse to the global-maximum (tallest) correlation peak.\\
5  & Arcane Capacitor      & $f\approx105.6$\,MHz.\\
6  & Counter-Flow HX       & $\dot m_h\approx0.80$\,kg/s, $\dot m_c\approx2.50$\,kg/s $\Rightarrow T_{h,out}\approx70\,^{\circ}$C.\\
7  & Sentry-Bot FK         & $\theta_1\approx16.8^{\circ}$, $\theta_2\approx17.7^{\circ}$, $\theta_3\approx-52.5^{\circ}$.\\
8  & Rankine Cycle         & $T_H\approx350\,^{\circ}$C, $T_L\approx30\,^{\circ}$C $\Rightarrow\eta\approx45\,\%$.\\
9  & Nyquist Hack          & $N=0$ (stable; $Z=N+P=0$).\\
10 & Torsional Stress      & $T_{max}=30.40$\,kN$\cdot$m.\\
11 & Void Frequency Tuner  & $L=3.57$\,mH $\Rightarrow f_0=2663$\,Hz.\\
12 & Beam Deflection       & Marker at $x\approx3.06$\,m.\\
13 & Singularity Trap      & $\theta_2\approx87^{\circ}$ (keep $14.5^{\circ}<\theta_2<165.5^{\circ}$).\\
14 & PID Stabilizer        & $K_p\approx0.40$, $K_i\approx0.05$, $K_d\approx0.55$.\\
15 & Mach Shock            & Shock at $x\approx0.81$ ($A_s/A^{*}\approx1.45$).\\
16 & Kalman Tracker        & Click blip-cluster centroid ($\approx x=35$, Round 1).\\
17 & Spectral De-Scrambler & Peaks at $45$, $212$, $395$\,Hz.\\
18 & Prism Shield (LQR)    & Balanced $\mathbf Q/\mathbf R\approx1$, centre of plot (slight high-$\mathbf Q$).\\
\bottomrule
\end{longtable}

\vspace{6pt}
\textit{Notes.} Tasks 3, 4, 16 and 18 depend on data that is hidden or
graphical in the live page (offline scanner, correlation trace, sensor blips,
empty selector); for these the governing method is given so the exact on-screen
value can be applied directly. All numeric answers were verified against the
live readouts shown in the task screenshots.

\end{document}